\chapter{Revisão de Literatura}\label{chap:revisao}

Este capítulo apresenta os fundamentos teóricos e os trabalhos relacionados
que embasam o desenvolvimento de um sistema computacional para interpretação
de análises de solo e recomendação agronômica. A revisão foi conduzida
de forma narrativa, com buscas realizadas nas bases Google Scholar, SciELO,
IEEE Xplore e no repositório da Embrapa, utilizando os termos
\textit{análise de solo}, \textit{calagem e adubação}, \textit{aptidão agrícola},
\textit{expert system agriculture} e \textit{sistema especialista agronomia},
com filtros para publicações em língua portuguesa e inglesa. Foram priorizadas
publicações com abrangência para o estado do Rio Grande do Sul e trabalhos
que descrevem sistemas computacionais aplicados ao manejo da fertilidade do solo.

\section{Fundamentação Teórica}

Esta seção apresenta os conceitos teóricos que sustentam o desenvolvimento
do trabalho. São abordados os principais parâmetros das análises químicas e
físicas de solo, os critérios de calagem e adubação vigentes para o Rio Grande
do Sul, os fundamentos da avaliação de aptidão agrícola das terras e os
princípios dos sistemas especialistas baseados em regras.

\subsection{Análise Química e Física do Solo}

A análise de solo é o principal instrumento para o diagnóstico da fertilidade
e para o planejamento das práticas de manejo nutricional das culturas. Por meio
dela, determinam-se os teores de nutrientes disponíveis, a reação do solo e
atributos físicos que condicionam a produtividade das lavouras
\citep{Embrapa2011}.

Os principais parâmetros avaliados em uma análise química de solo incluem o
pH em água, que indica a acidez ou alcalinidade do meio e condiciona a
disponibilidade da maioria dos nutrientes; o teor de fósforo disponível (P);
os teores de potássio (K), cálcio (Ca) e magnésio (Mg) trocáveis; o alumínio
trocável (Al\textsuperscript{3+}), associado à toxidez em solos ácidos;
a capacidade de troca de cátions (CTC), que expressa a capacidade do solo
em reter e fornecer nutrientes; e a saturação por bases (V\%), que indica
a proporção da CTC ocupada por bases trocáveis \citep{Embrapa2011, ManualCalagem2016}.

A textura do solo, determinada pela proporção de areia, silte e argila,
constitui o principal atributo físico considerado na interpretação das análises,
pois influencia diretamente a CTC, a capacidade de retenção de água e a
resposta do solo à calagem \citep{ManualCalagem2016}. O teor de matéria orgânica
(MO) também é avaliado, dada sua contribuição para a CTC efetiva e para a
disponibilidade de nitrogênio \citep{Santos2003}. Além dos macronutrientes,
os micronutrientes --- como boro, manganês, cobre e zinco --- vêm recebendo
atenção crescente nas pesquisas regionais, com a publicação recente de
conjuntos de dados experimentais sobre sua aplicação em culturas de grãos
no Sul do Brasil \citep{FontouraEtAl2026}.

Conforme \cite{EmbrapaDoc206}, a interpretação dos resultados analíticos exige
que os valores obtidos sejam comparados a faixas de referência estabelecidas
para cada parâmetro, considerando-se o tipo de solo e a cultura de interesse.

\subsection{Calagem e Adubação para o Rio Grande do Sul}

O Manual de Calagem e Adubação para os estados do Rio Grande do Sul e de
Santa Catarina, elaborado pela Comissão de Química e Fertilidade do Solo
(CQFS/NRS-SBCS), constitui a principal referência técnica para as
recomendações de correção e fertilização no sul do Brasil
\citep{ManualCalagem2016}. O manual estabelece critérios específicos para
interpretação dos atributos do solo, métodos de cálculo da necessidade de
calagem e doses de fertilizantes para as principais culturas da região.

A calagem é a prática que visa elevar o pH do solo a níveis adequados para
o desenvolvimento das culturas, neutralizando o alumínio tóxico e elevando
a disponibilidade de cálcio e magnésio. Para a maioria das culturas de interesse
regional, o manual define um pH de referência específico, a partir do qual
o cálculo da necessidade de calagem é realizado com base na saturação por bases
atual do solo, na saturação alvo correspondente ao pH de referência da cultura,
na CTC e no poder de neutralização do corretivo \citep{ManualCalagem2016}.

A adubação fosfatada e potássica é dimensionada com base nos teores
disponíveis no solo, interpretados em classes de disponibilidade (muito baixo,
baixo, médio, alto e muito alto), e na expectativa de rendimento da cultura.
O manual fornece tabelas de recomendação individuais para cada cultura ou grupo
de culturas \citep{ManualCalagem2016}.

\subsection{Aptidão Agrícola das Terras}

A avaliação da aptidão agrícola das terras tem por objetivo classificar
as áreas de acordo com seu potencial para os diferentes tipos de uso
agropecuário, considerando as limitações impostas pelos atributos do solo
e do relevo sob diferentes níveis tecnológicos. O principal sistema utilizado
no Brasil é o Sistema de Avaliação da Aptidão Agrícola das Terras, desenvolvido
pela Embrapa, que classifica as terras em seis grupos de aptidão com base em
cinco fatores de limitação: deficiência de fertilidade, deficiência hídrica,
excesso de água, suscetibilidade à erosão e impedimentos à mecanização
\citep{EmbrapaAptidao2025}.

Conforme apontado por \cite{Hofig2015}, em estudo conduzido no município de
Cerro Grande do Sul/RS, o Sistema de Avaliação da Aptidão Agrícola das Terras
é mais adequado para análises em âmbito regional, não indicando práticas de
manejo para os diferentes tipos de utilização nem sendo projetado para a
avaliação individualizada de propriedades a partir de dados de análise de solo.
Os autores ressaltam que adaptações são necessárias para seu uso no planejamento
conservacionista de propriedades rurais, evidenciando uma lacuna metodológica
relevante para propriedades rurais de escala individual.

Uma abordagem alternativa consiste na estimativa de aptidão a partir dos
próprios parâmetros da análise de solo, cruzados com as exigências agronômicas
de cada cultura, sem considerar fatores de relevo e clima. Trata-se de uma
estimativa de aptidão edáfica, delimitada às condições de fertilidade e físicas
do solo para uma cultura específica, mais adequada à avaliação individualizada
de áreas a partir de laudos de análise.

\section{Sistemas Especialistas Baseados em Regras}

Sistemas especialistas (SE) são programas computacionais que utilizam
técnicas de inteligência artificial para reproduzir o raciocínio de um
especialista humano em um domínio específico do conhecimento, apoiando
processos de tomada de decisão, diagnóstico e recomendação
\citep{SantosEtAl1994}. Conforme \cite{SantosEtAl1994}, um sistema
especialista é constituído por três componentes fundamentais: a
\textbf{base de conhecimento}, que armazena os fatos, regras e heurísticas
do domínio; o \textbf{motor de inferência}, responsável por aplicar as
regras à base de conhecimento para derivar conclusões a partir dos dados
de entrada; e a \textbf{interface com o usuário}, que permite a comunicação
bidirecional entre o sistema e o usuário final.

Os sistemas baseados em regras do tipo SE-ENTÃO (\textit{if-then})
constituem a abordagem mais difundida para a representação do conhecimento
agronômico em sistemas computacionais. Nesse modelo, cada regra representa
um fragmento do conhecimento especialista na forma
\textit{SE} <condição> \textit{ENTÃO} <conclusão>, e o motor de inferência
percorre a base de regras para identificar quais condições são satisfeitas
pelos dados de entrada, derivando as recomendações correspondentes
\citep{Lanzer2010}.

A aplicação de sistemas especialistas na agricultura tem sido estudada em
diferentes contextos. \cite{Lanzer2010} desenvolveu um sistema especialista
baseado em regras para a classificação do uso e manejo das terras, aplicando
representação hierárquica do conhecimento~--- indicada para sistemas que lidam
com classificação~--- em um sistema de determinação de uso das terras compatível
com o sistema brasileiro. O autor destaca que a introspecção facilita o
desenvolvimento nos estágios iniciais de um sistema especialista quando as
informações para definição dos estados estão presentes nos próprios sistemas
de classificação, como é o caso do Manual de Calagem e Adubação
\citep{ManualCalagem2016}.

Em escala mais ampla, \cite{Shahzadi2016} propõem um sistema especialista
para agricultura inteligente baseado em Internet das Coisas (IoT), no qual
regras de conhecimento são alimentadas com dados coletados em tempo real por
sensores, gerando diagnósticos e recomendações preventivas para o manejo de
doenças e pragas. Embora voltado a um domínio diferente, o estudo evidencia
a viabilidade e a aplicabilidade prática de sistemas especialistas baseados
em regras para o suporte à decisão agronômica.

\section{Trabalhos Relacionados}

Esta seção apresenta os principais sistemas e ferramentas computacionais
identificados na literatura que se relacionam ao escopo do presente trabalho,
discutindo suas características, abrangência e limitações.

\textbf{FertFacil} é um sistema online e gratuito desenvolvido para estimar
a necessidade de adubação e calagem para os principais cultivos do RS, SC,
MS e Paraguai \citep{Tiecher2016}. O sistema foi desenvolvido com base nos
critérios do Manual de Calagem e Adubação e permite ao técnico ou produtor
obter recomendações de forma rápida, a partir da entrada dos resultados da
análise de solo e da cultura desejada. Atualmente o FertFacil cobre mais de
100 cultivos e é a ferramenta de escolha da Emater/RS, sendo utilizado em
mais de 480 municípios do estado. Sua principal limitação identificada é a
ausência de um módulo de estimativa de aptidão agrícola integrado ao processo
de recomendação.

\textbf{AdubaTec} é um sistema web gratuito desenvolvido pela Embrapa
Mandioca e Fruticultura em parceria com a Universidade Federal do
Recôncavo da Bahia (UFRB), voltado à recomendação de calagem e adubação
para mandioca e diversas fruteiras tropicais, com arquitetura parametrizável
que permite a incorporação de novas culturas sem modificação do código-fonte
\citep{AdubaTec2021}. O AdubaTec é baseado em manuais de recomendação
voltados ao Nordeste do Brasil, não cobrindo os critérios do Manual de
Calagem e Adubação do RS/SC, e não contempla a estimativa de aptidão
agrícola a partir da análise individual de solo.

\textbf{Sistema de Avaliação de Aptidão Agrícola das Terras da Embrapa}
\citep{EmbrapaAptidao2025} é o principal sistema de referência para
classificação do potencial produtivo das terras brasileiras. Trata-se de
um sistema voltado a levantamentos regionais em escala exploratória, baseado
em mapeamento pedológico e em fatores de limitação que incluem relevo,
clima e solo. Sua escala de análise o torna inadequado para a avaliação
individualizada de uma propriedade a partir de dados de análise de solo,
conforme discutido por \cite{Hofig2015}.

A partir do levantamento realizado, constata-se que os sistemas disponíveis
abordam, de forma separada, a recomendação de manejo do solo ou a avaliação
de aptidão agrícola em escala regional. Nenhum deles integra ambas as
perspectivas a partir de uma análise individual de solo com cobertura para
o Manual RS/SC, evidenciando uma lacuna que justifica o desenvolvimento de
uma solução específica para esse contexto.
